% Generated draft: requires mathematical and source review.
\documentclass[11pt]{article}
\usepackage[margin=25mm]{geometry}
\usepackage{fontspec}
\setmainfont{Latin Modern Roman}
\usepackage{amsmath,amssymb,mathtools,unicode-math}
\setmathfont{Latin Modern Math}
\usepackage{enumitem,needspace,xurl,hyperref}
\setlength{\parindent}{0pt}
\setlength{\parskip}{5pt}
\setlength{\emergencystretch}{3em}
\newcommand{\N}{\mathbb{N}}
\newcommand{\Z}{\mathbb{Z}}
\newcommand{\Q}{\mathbb{Q}}
\newcommand{\R}{\mathbb{R}}
\newcommand{\C}{\mathbb{C}}
\newcommand{\F}{\mathbb{F}}
\newcommand{\A}{\mathbb{A}}
\newcommand{\PP}{\mathbb{P}}
\newcommand{\E}{\mathbb{E}}
\newcommand{\eps}{\varepsilon}
\newcommand{\dd}{\,\mathrm{d}}
\begin{document}
\section*{UnsolvedMath ID: 30002128}
\textbf{Source problem code:} catalogue entry
\subsection{Definitions and mathematical statement}
Let \(F=\langle a,b\rangle\) be the free group of rank two. The supplied record does not specify the target of the representations or the choice of lifts. To make the notation explicit, consider the conventional lifted circle-action interpretation
\[
\widetilde{\operatorname{Homeo}}^{+}(S^{1})
=\{f:\R\to\R\mid f\text{ is an increasing homeomorphism and }f(x+1)=f(x)+1\text{ for every }x\in\R\},
\qquad S^{1}=\R/\Z.
\]
For an element \(f\) of this group, define its lifted rotation number, also called its translation number, by
\[
\tau(f)=\lim_{n\to\infty}\frac{f^{n}(x)-x}{n},
\]
where the limit is independent of \(x\in\R\). Its reduction modulo \(\Z\) is the rotation number of the induced circle homeomorphism. Under this interpretation, for \(w\in F\) and \(r,s\in\R\), set
\[
\mathcal H(r,s)=\left\{\rho:F\longrightarrow\widetilde{\operatorname{Homeo}}^{+}(S^{1})\ \middle|\ \rho\text{ is a group homomorphism},\ \tau(\rho(a))=r,\ \tau(\rho(b))=s\right\},
\]
\[
R(w,r,s)=\sup_{\rho\in\mathcal H(r,s)}\tau(\rho(w)).
\]
This is an explicit conventional interpretation, not a representation target independently confirmed from the cited report.

The supplied question asks whether
\[
\text{for every }w\in F\text{ and every }r,s\in\R,\qquad R(w,r,s)\in\Q,
\]
and asks for a practical method to compute or approximate \(R(w,r,s)\).

\textbf{Unresolved formulation qualification.} Under the interpretation above, the generator constraints themselves give \(R(a,r,s)=r\) and \(R(b,r,s)=s\). Thus the universal rationality clause, which includes generator words and arbitrary irrational inputs, is incompatible with these definitions. The intended restrictions on the words or inputs, or the intended alternative rationality assertion, must be checked against the original source. No replacement conjecture is selected here. The computational request also leaves unspecified how arbitrary real inputs are represented, what constitutes a practical method, and what approximation error or guarantee is required.
\subsection{Short English statement}
The supplied problem concerns the largest lifted rotation number obtainable for a word in two free generators when their lifted rotation numbers are prescribed. It asks about rationality and practical computation or approximation. The representation target and lift conventions require confirmation, and the stated rationality claim for all real inputs conflicts with the generator cases under the conventional interpretation.
\subsection{Sources}
\begin{itemize}
\item Supplied local UnsolvedMath registry record: permanent ID \(30002128\), exact source problem code \texttt{catalogue entry}, title ``Rationality and Computation of Free-Group Rotation Suprema,'' in the Oberwolfach Reports Open Problems collection. The record assigns proposed year \(2012\); no individual proposer is identified.
\item Supplied report citation: \emph{Geometric Group Theory, Hyperbolic Dynamics and Symplectic Geometry (2013)}, Oberwolfach Reports, DOI \texttt{10.4171/owr/2012/35}, pp.\ 2139--2203. \url{https://doi.org/10.4171/owr/2012/35}. The supplied background describes Oberwolfach Reports as reports of mathematical workshops and their open-problem discussions. The citation's parenthetical year is retained despite the record's proposed year and the DOI's year component being \(2012\).
\item Supplied literature source: Calegari--Walker, \emph{Ziggurats and rotation numbers}. \url{https://arxiv.org/abs/1110.0080}.
\item Supplied literature source: Mann, \emph{rigidity and flexibility of circle actions}. \url{https://e.math.cornell.edu/people/mann/papers/rigflexA.pdf}.
\item Supplied literature source: \emph{Oberwolfach Report 35}. \url{https://doi.org/10.4171/owr/2012/35}.
\item Literature assessment dated exactly \texttt{2026-08-21}: the supplied curation labels the problem open and states that rationality of the supremal rotation function for every word and all real rotation inputs is not proved, that primary work establishes important special cases, and that the general problem remains open. This is supplied dated curation derived from the OWR record, not an independently verified assessment. Its formulation requires review in light of the generator-input conflict described above; time-sensitive status and cited claims require independent verification.
\item The registry labels the statement \texttt{corrected\_verified} and supplies a verification note claiming comparison with the official Oberwolfach report and normalization of notation, context, and question form without changing the source-level claim. These are supplied provenance claims, not verification performed for this draft. No linked source was visited in preparing this statement.
\end{itemize}
\end{document}
